\chapter{Formula Index}
\label{app:formulas}
Every formula below is derived somewhere in the book; the last column says where, so
this index doubles as a map back to the reasoning. Nothing here needs to be memorized
that can be re-derived from the element laws, KCL, and KVL.

\begin{longtable}{@{}L{0.31\textwidth}L{0.47\textwidth}L{0.14\textwidth}@{}}
\toprule
Formula & Use & Derived in \\
\midrule
\endhead
\multicolumn{3}{@{}l}{\textbf{Elements and energy}}\\
\(i=\dd q/\dd t\) & Current as rate of charge flow. & \cref{ch:foundations} \\
\(V=IR\) & Resistor law. & \cref{ch:foundations} \\
\(P=VI=I^2R=V^2/R\) & Power in resistive circuits. & \cref{ch:resistive} \\
\(i=C\,\dd v/\dd t\) & Capacitor current-voltage law. & \cref{ch:foundations} \\
\(v=L\,\dd i/\dd t\) & Inductor voltage-current law. & \cref{ch:foundations} \\
\(E_C=\frac{1}{2}CV^2\) & Capacitor energy. & \cref{ch:foundations} \\
\(E_L=\frac{1}{2}LI^2\) & Inductor energy. & \cref{ch:foundations} \\
\(f_{\mathrm{SRF}}=1/(2\pi\sqrt{L_sC})\) & Self-resonance of a real part; above it the part inverts type. & \cref{sec:selfresonance} \\
\(\delta=\sqrt{2/(\omega\mu\sigma)}=1/\sqrt{\pi f\mu\sigma}\) & Skin depth, from the diffusion equation; \SI{66}{\micro\meter} in copper at \SI{1}{\mega\hertz}. & \cref{sec:selfresonance} \\
\(R_{\mathrm{AC}}\propto1/\delta\propto\sqrt f\) & Skin effect raises AC resistance; \(R_{\mathrm{AC}}/R_{\mathrm{DC}}\approx a/2\delta\). & \cref{sec:selfresonance} \\
\(f_{\text{peak}}=f_{\mathrm{SRF}}/\sqrt5\) & Where a coil's \(Q\) peaks: \(\sqrt f\) rise against the \(C_p\) fall. & \cref{sec:selfresonance} \\
\(R_{\mathrm{eq}}=\sum R_k\), \(1/C_{\mathrm{eq}}=\sum 1/C_k\) & Series combining (parallel is the dual). & \cref{ch:resistive} \\
\addlinespace
\multicolumn{3}{@{}l}{\textbf{Decibels}}\\
\(10\log_{10}(P_2/P_1)\) & Power ratio in dB. & \cref{sec:decibels} \\
\(20\log_{10}(V_2/V_1)\) & Amplitude ratio (same impedance). & \cref{sec:decibels} \\
\(-\SI{3.01}{\decibel}=20\log_{10}(1/\sqrt2)\) & Half power; the cutoff convention. & \cref{sec:decibels} \\
\addlinespace
\multicolumn{3}{@{}l}{\textbf{First-order circuits}}\\
\(\tau_{RC}=RC\) & RC time constant. & \cref{ch:rc} \\
\(\tau_{RL}=L/R\) & RL time constant. & \cref{ch:rl} \\
\(G_C=1/(1+RCs)\) & RC low-pass (output across \(C\)). & \cref{sec:rc-freq} \\
\(G_R=RCs/(1+RCs)\) & RC high-pass (output across \(R\)). & \cref{sec:rc-freq} \\
\addlinespace
\multicolumn{3}{@{}l}{\textbf{Signals and spectra}}\\
\(x=\sum_n c_n\ee^{\jj n\omega_0t}\) & Fourier series of a periodic signal. & \cref{sec:fourier} \\
\(y=\sum_n G(\jj n\omega_0)c_n\ee^{\jj n\omega_0 t}\) & \(G\) at the harmonics is the whole answer---the licence for analyzing one sinusoid at a time. & \cref{sec:fourier} \\
\(x(t+T/2)=-x(t)\) & Half-wave symmetry forbids even harmonics (square wave). & \cref{sec:fourier} \\
\addlinespace
\multicolumn{3}{@{}l}{\textbf{Higher order: cascades, pole patterns, delay}}\\
\(G=K\prod\frac{1}{1+\tau_ks}\prod\frac{\omega_{0j}^2}{s^2+2\zeta_j\omega_{0j}s+\omega_{0j}^2}\) & Factoring theorem: every real \(G(s)\) is first- and second-order sections. & \cref{sec:factoring} \\
\(20\log|G|=\sum_k 20\log|G_k|\) & Cascades add in decibels; phases add in degrees. & \cref{sec:cascadeadd} \\
\(-20n\ \si{\decibel}\)/decade, \(-90n^\circ\) & Roll-off and phase lag of an \(n\)-pole response. & \cref{sec:cascadeadd} \\
\(\omega_{\mathrm{osc}}=\sqrt3\,\omega_0\) at \(K=8\) & Three identical lags: the gain and frequency at which a loop oscillates. & \cref{sec:threepoles} \\
\(G(s)=\ee^{-sT}\) & Pure delay: \(|G|=1\), \(\angle G=-\omega T\) without bound. & \cref{sec:puredelay} \\
\(|G|^2=1/[1+(\omega/\omega_c)^{2n}]\) & Butterworth magnitude; poles equally spaced on a semicircle. & \cref{sec:poleplacement} \\
\(Q=1/(2\cos\phi)\) & Section \(Q\) of a pole pair at angle \(\phi\); source of the Butterworth table. & \cref{sec:poleplacement} \\
\(\tau_g=-\dd\phi/\dd\omega\) & Group delay; constant \(\tau_g\) is linear phase. & \cref{sec:groupdelay} \\
\addlinespace
\multicolumn{3}{@{}l}{\textbf{Phasors, impedance, power}}\\
\(Z_R=R\) & Resistor impedance. & \cref{ch:ac} \\
\(Z_L=\jj\omega L\) & Inductor impedance. & \cref{ch:ac} \\
\(Z_C=1/(\jj\omega C)\) & Capacitor impedance. & \cref{ch:ac} \\
\(\omega=2\pi f\) & Angular frequency conversion. & \cref{ch:foundations} \\
\(S=VI^*\) & Complex power using RMS phasors. & \cref{ch:ac} \\
\addlinespace
\multicolumn{3}{@{}l}{\textbf{RMS, average power, and PEP}}\\
\(X_{\RMS}=\sqrt{\frac1T\int_0^T x^2\dd t}\) & Definition of RMS (any waveform). & \cref{sec:rms} \\
\(X_{\RMS}^2=X_{\mathrm{dc}}^2+\sum_n X_n^2\) & Parseval: RMS adds in quadrature over harmonics. & \cref{sec:parseval} \\
\(\pi/(2\sqrt2)\approx1.111\) & Form factor of a sine; why averaging meters misread. & \cref{sec:parseval} \\
\(V_{\RMS}=V_p/\sqrt2=V_{pp}/(2\sqrt2)\) & Sinusoid only; \(V_{pp}=2V_p\). & \cref{sec:rms} \\
\(P=V_{\RMS}^2/R=I_{\RMS}^2R\) & Average power; DC formulas hold with RMS values. & \cref{sec:pep} \\
\(P=V_{pp}^2/(8R)\) & Power straight from a scope reading. & \cref{sec:pep} \\
\(\mathrm{PEP}=V_{pp,\max}^2/(8R)\) & Peak envelope power; equals average power for an unmodulated carrier. & \cref{sec:pep} \\
\addlinespace
\multicolumn{3}{@{}l}{\textbf{Resonance and second-order systems}}\\
\(\omega_0=1/\sqrt{LC}\) & Undamped LC angular frequency. & \cref{ch:rlc} \\
\(f_0=1/(2\pi\sqrt{LC})\) & Resonant frequency in hertz. & \cref{ch:rlc} \\
\(s^2+2\zeta\omega_0 s+\omega_0^2\) & Standard second-order pole polynomial. & \cref{ch:linsys} \\
\(\det(sI-A)=0\) & Poles are the eigenvalues of \(A\). & \cref{ch:rlc} \\
\(\alpha=R/2L,\ \zeta=\alpha/\omega_0\) & Series RLC damping rate and ratio. & \cref{ch:rlc} \\
\(s=-\alpha\pm\jj\omega_d,\ \omega_d=\sqrt{\omega_0^2-\alpha^2}\) & Underdamped pole pair. & \cref{ch:rlc} \\
\(Q=1/(2\zeta)\) & Quality factor and damping ratio. & \cref{ch:linsys} \\
\(Q_s=\omega_0L/R\) & \emph{Series} RLC quality factor. & \cref{ch:rlc} \\
\(Q_p=R/(\omega_0L)=\omega_0RC\) & \emph{Parallel} RLC quality factor. & \cref{ch:rlcpar} \\
\(R_p=(1+Q^2)R_s\), \(X_p=(1+Q^{-2})X_s\) & Series\(\leftrightarrow\)parallel at one frequency; \(Q\) is unchanged. & \cref{sec:seriesparallel} \\
\(Z_{\max}=(1+Q^2)r\approx L/(Cr)\) & Real tank's peak impedance, from its coil loss \(r\). & \cref{sec:seriesparallel} \\
\(BW=f_0/Q\) & Half-power bandwidth; exact, not an estimate. & \cref{sec:halfpower} \\
\(f_0=\sqrt{f_1f_2}\) & Half-power edges straddle resonance geometrically. & \cref{sec:halfpower} \\
\(G=\omega_0^2/(s^2+\tfrac{\omega_0}{Q}s+\omega_0^2)\) & LC low-pass as a filter. & \cref{sec:lc-lowpass} \\
\addlinespace
\multicolumn{3}{@{}l}{\textbf{Feedback and active circuits}}\\
\(A_{\mathrm{cl}}=A/(1+A\beta)\) & Closed-loop gain; \(L=A\beta\) is loop gain. & \cref{ch:feedback} \\
\(A_{\mathrm{cl}}\to1/\beta\) & Large-loop-gain limit. & \cref{ch:feedback} \\
\(\dd A_{\mathrm{cl}}/A_{\mathrm{cl}}=(\dd A/A)/(1+L)\) & Feedback divides gain error by \(1+L\). & \cref{ch:feedback} \\
\(Z_{\mathrm{out}}=R_o/(1+L)\), \(Z_{\mathrm{in}}=R_{id}(1+L)\) & Feedback sets impedance; the op-amp ideals are the \(L\to\infty\) limit. & \cref{sec:feedbackz} \\
\(\text{gain}\times\text{BW}=A_0\omega_a\) & Gain--bandwidth product (single-pole amp). & \cref{sec:gbw} \\
\(\beta=R_g/(R_f+R_g)\) & Feedback fraction from the resistors. & \cref{sec:beta-resistors} \\
\(-R_f/R_{in}\), \(1+R_f/R_g\) & Inverting and non-inverting gains. & \cref{ch:active} \\
\(L=-1\): \(|L|=1,\ \angle L=-180^\circ\) & Barkhausen oscillation condition. & \cref{ch:feedback} \\
\(\omega_0=1/\sqrt{R_1R_2C_1C_2}\) & Sallen--Key corner frequency. & \cref{ch:active} \\
\(f_p=f_s\sqrt{1+C_m/C_0}\) & Crystal's two resonances; the gap \(\approx C_m/2C_0\) is the pulling range. & \cref{sec:crystal} \\
\(L(s)=K_dK_vF(s)/s\) & PLL loop gain; the VCO supplies the free integrator. & \cref{sec:pll} \\
\(\phi_e(\infty)=\Delta\omega/K\) & Static phase error of a type-1 loop (frequency locks exactly). & \cref{sec:pll} \\
\(\omega_n=\sqrt{K\omega_f},\ \zeta=\tfrac12\sqrt{\omega_f/K}\) & PLL dynamics with a single-pole loop filter. & \cref{sec:pll} \\
\(f_{\mathrm{out}}=Nf_{\mathrm{ref}}\) & Synthesizer; reference noise rises \(20\log_{10}N\). & \cref{sec:pll} \\
\(\eta=\tfrac12(I_1/I_{\mathrm{dc}})(V_1/V_{\mathrm{dc}})\) & Amplifier efficiency: conduction angle times load coupling. & \cref{sec:conductionangle} \\
\(\eta=\tfrac12\) (Class A), \(\pi/4\) (Class B) & From the truncated-cosine Fourier coefficients. & \cref{sec:conductionangle} \\
\(p=vi=0\) for an ideal switch & Why switching amplifiers have no efficiency ceiling. & \cref{sec:switching} \\
\(f(x)-f(-x)=2\sum_{n\ \mathrm{odd}}a_nx^n\) & Push-pull: even orders vanish for \emph{any} \(f\); matching is the requirement. & \cref{sec:pushpull} \\
\(f_{\mathrm{LO}}\pm f_{\mathrm{RF}}\) & Mixer output frequencies (nonlinear!). & \cref{sec:mixers} \\
\(2f_{\mathrm{IF}}\) & Image offset from the wanted signal. & \cref{sec:mixers} \\
\(2f_1-f_2,\ 2f_2-f_1\) & Third-order intermodulation products; the two-tone test. & \cref{sec:mixers} \\
\(C_{\text{neut}}=C_{\text{stray}}\) & Neutralization: cancel the device's own feedback path. & \cref{sec:neutralization} \\
\(2f_{\text{line}}\) & Ripple frequency of a full-wave rectifier (\(f_{\text{line}}\) if half-wave). & \cref{ch:active} \\
\addlinespace
\multicolumn{3}{@{}l}{\textbf{Matching and transformers}}\\
\(Q=\sqrt{R_{\text{hi}}/R_{\text{lo}}-1}\) & L-network matching \(Q\); solves \(R_{\text{hi}}=R_{\text{lo}}(1+Q^2)\). & \cref{sec:lnetwork} \\
\(X_{\text{series}}=QR_{\text{lo}}\), \(X_{\text{shunt}}=R_{\text{hi}}/Q\) & L-network element reactances. & \cref{sec:lnetwork} \\
\(Z_p=(N_p/N_s)^2Z_s\) & Transformer impedance ratio. & \cref{ch:filters} \\
\addlinespace
\multicolumn{3}{@{}l}{\textbf{Transmission lines}}\\
\(Z_0=\sqrt{L'/C'},\ v_p=1/\sqrt{L'C'}\) & Line impedance and phase velocity. & \cref{ch:lines} \\
\(\lambda=v_p/f=VF\,c/f\) & Wavelength on the line. & \cref{ch:lines} \\
\(\Gamma=(Z_L-Z_0)/(Z_L+Z_0)\) & Load reflection coefficient. & \cref{ch:lines} \\
\(a=(V+Z_0I)/2\sqrt{Z_0}\), \(b=(V-Z_0I)/2\sqrt{Z_0}\) & Wave coordinates; \(|a|^2\) is incident power. & \cref{sec:sparams} \\
\(S_{11}=b/a=\Gamma\) & One-port \(S\)-parameter \emph{is} the reflection coefficient. & \cref{sec:sparams} \\
\(S_{21}=b_2/a_1\) & Forward transmission; a sweep of it is a Bode plot of \(G(\jj\omega)\). & \cref{sec:sparams} \\
\(V_2/V_1=S_{21}/(1+S_{11})\) & \(S_{21}\) equals the voltage gain only when the input is matched. & \cref{sec:sparams} \\
\(\mathrm{SWR}=(1+|\Gamma|)/(1-|\Gamma|)\) & Standing wave ratio. & \cref{ch:lines} \\
\(\mathrm{RL}=-20\log_{10}|\Gamma|\) & Return loss in decibels. & \cref{ch:lines} \\
\(Z_{\mathrm{in}}=Z_0\dfrac{Z_L+\jj Z_0\tan\beta\ell}{Z_0+\jj Z_L\tan\beta\ell}\) & Input impedance along a lossless line. & \cref{ch:lines} \\
\(Z_{\mathrm{in}}=Z_0^2/Z_L\) & Quarter-wave impedance inversion. & \cref{ch:lines} \\
\(Z_t=\sqrt{Z_0Z_L}\) & Quarter-wave transformer for real impedances. & \cref{ch:lines} \\
\(\ell_{\lambda/2}\approx468/f_{\mathrm{MHz}}\) ft & Resonant dipole length (end effect included); \(234/f\) for a quarter wave. & \cref{sec:antennalength} \\
\(\eta=R_{\mathrm{rad}}/(R_{\mathrm{rad}}+R_{\mathrm{loss}})\) & Antenna efficiency; loss \emph{improves} SWR while lowering it. & \cref{sec:antennaeff} \\
\(R(z)=R_{\mathrm{centre}}/\cos^2\beta z\) & Feed-point resistance vs position; \(2\times\) at \(\lambda/8\), \(\to\infty\) at the end. & \cref{sec:feedpoint} \\
\(|F|=2|\cos((\phi+\beta d\cos\theta)/2)|\) & Two-element array factor; the four named patterns are substitutions. & \cref{sec:arrayfactor} \\
\(\mathrm{dBi}=\mathrm{dBd}+2.15\) & Antenna gain references: isotropic vs dipole. & \cref{sec:antennalength} \\
\(\jj Z_0\tan\beta\ell\) & Shorted-stub reactance. & \cref{sec:stubs} \\
\(\gamma=\sqrt{(R'+\jj\omega L')(G'+\jj\omega C')}=\alpha+\jj\beta\) & Complex propagation constant of a real line. & \cref{sec:lineloss} \\
\(\alpha\approx R'/(2Z_0)+G'Z_0/2\) & Low-loss attenuation constant: conductor plus dielectric. & \cref{sec:lineloss} \\
\(\text{dB}\approx8.686\,\alpha x\) & Nepers to decibels; loss is proportional to length. & \cref{sec:lineloss} \\
\(|\Gamma_{\mathrm{in}}|=|\Gamma_L|\ee^{-2\alpha\ell}\) & Line loss masks a mismatch (round trip). & \cref{sec:lineloss} \\
\addlinespace
\multicolumn{3}{@{}l}{\textbf{Sampling and digital signal processing}}\\
\(X_s(f)=\frac{1}{T_s}\sum_n X(f-nf_s)\) & Sampling replicates the spectrum at every multiple of \(f_s\). & \cref{sec:samplingismult} \\
\(f_s>2f_{\max}\) & Sampling theorem; \(f_s/2\) is the Nyquist frequency. & \cref{sec:samplingthm} \\
\(f_{\text{alias}}=|f-nf_s|\) & Where an out-of-band component folds to. & \cref{sec:samplingthm} \\
\(\mathrm{SNR}_{\dB}=6.02N+1.76\) & Quantization noise: \SI{6}{\decibel} per bit. & \cref{sec:quantnoise} \\
\(10\log_{10}\bigl(\tfrac{f_s/2}{B}\bigr)\) & Oversampling processing gain. & \cref{sec:quantnoise} \\
\(\tau_g=\tfrac{N-1}{2}T_s\) & Symmetric FIR: constant group delay, exactly linear phase. & \cref{sec:firiir} \\
\(f_{\mathrm{out}}=f_{\mathrm{clk}}M/2^N\) & Direct digital synthesis output frequency. & \cref{ch:dsp} \\
\addlinespace
\multicolumn{3}{@{}l}{\textbf{Noise, sensitivity, and dynamic range}}\\
\(P_n=kTB\) & Available thermal noise power; independent of \(R\). & \cref{sec:noisefloor} \\
\(-174+10\log_{10}B+\mathrm{NF}\) & Receiver noise floor in dBm; also the MDS. & \cref{sec:noisefloor} \\
\(F=1+T_e/T_0\) & Noise figure and equivalent noise temperature. & \cref{sec:noisefigure} \\
\(F=F_1+\frac{F_2-1}{G_1}+\cdots\) & Friis: the first stage dominates. & \cref{sec:friis} \\
\(y=a_1x+a_2x^2+a_3x^3\) & Weak nonlinearity; slopes 2 and 3 dB/dB follow. & \cref{sec:powerseries} \\
\(\mathrm{IIP3}\approx P_{\SI{1}{\decibel}}+\SI{9.6}{\decibel}\) & Cubic nonlinearity ties the two together. & \cref{sec:powerseries} \\
\(\mathrm{DR3}=\tfrac23(\mathrm{IIP3}-\mathrm{MDS})\) & IMD dynamic range; \(\tfrac{n-1}{n}\) for order \(n\). & \cref{sec:dynamicrange} \\
\(\mathrm{BDR}=P_{\SI{1}{\decibel}}-\mathrm{MDS}\) & Blocking dynamic range (a plain difference). & \cref{sec:dynamicrange} \\
\(\mathcal{L}\propto 1/(Q_L^2\Delta f^2)\) & Phase-noise skirt; doubling \(Q_L\) buys \SI{6}{\decibel}. & \cref{sec:phasenoise} \\
\(B_n=\int|H|^2\dd f/|H|^2_{\max}\) & Noise bandwidth; \(1.57f_c\) for one pole. & \cref{sec:noisebw} \\
\addlinespace
\multicolumn{3}{@{}l}{\textbf{Measurement}}\\
\(R_{\mathrm{in}}/(R_{\mathrm{in}}+R_{Th})\) & Voltmeter loading error. & \cref{ch:measurement} \\
\(R_1C_1=R_2C_2\) & Compensated scope probe (zero on pole). & \cref{sec:probe} \\
\bottomrule
\end{longtable}
